% HD-S30. GPT-6 Astra Pro, /root/sixfolds; inherited effort not exposed.
% 21 November 2026. Exact induced-block and polarized-period maps.
\subsection{HD-S30: rank-two induced blocks have period rank at most seven}
\label{sec:rank-two-blocks}

Let G be a finite group and let \(V\) be an absolutely irreducible
representation over an imaginary quadratic field K, with character
field exactly K. Suppose that
\[
                 V=W_1\oplus\cdots\oplus W_N
\]
is an actual G-permuted system of K-subspaces, each of dimension two.
For every connected G-cover of \(\mathbb P^1\) whose primitive
V-factor is a Weil sixfold, the period differential has rank at most
seven. This bound also holds after any specified K-linear isogeny
of that factor. In particular it applies to the natural block
representation of \(G_{12}\wr\mathfrak S_N\), with the literal
quadratic-field matrices of HD-S29 on each block.

The proof constructs the quotient coefficient cover and the exact
polarized variation through which the sixfold period map factors.
The stated hypothesis requires the block subspaces and their action;
it is not inferred from the dimension of V or from a group name.

\paragraph{The actual quotient tower and induced-module map.}
Transitivity on the blocks follows from irreducibility. Fix
\(W=W_1\), let \(H=\operatorname{Stab}_G(W)\), and let
\(H_0=\ker(H\to\operatorname{GL}(W))\). Put
\(\Gamma=H/H_0\). For the actual G-cover C define
\[
 C\longrightarrow X=C/H_0\xrightarrow{q}Y=C/H
                  \xrightarrow{h}\mathbb P^1.
\]
The map q is a connected Galois cover with finite deck group
\(\Gamma\), and \(\deg h=N\). All of these are quotient
morphisms of smooth compact curves. The group \(\Gamma\)
and its representation W are fixed along the marked G-Hurwitz
component. The H-module W is absolutely irreducible: an H-invariant
proper subspace would give the G-invariant proper direct sum of its
translates in the different blocks.

Choose an H-invariant positive Hermitian form on W by averaging.
Declare the blocks orthogonal and transport this form to every block;
this is well-defined and G-invariant. Write \(\pi_W\) for the
orthogonal block projection. The induced module is realized as
\[
 \mathcal I=\{f:G\to W:f(ax)=a f(x)\text{ for }a\in H\},
 \qquad (g f)(x)=f(xg).
\]
The following mutually inverse G-linear maps are literal:
\[
 V\longrightarrow\mathcal I,\quad
       v\longmapsto f_v(x)=\pi_W(xv),\qquad
 \mathcal I\longrightarrow V,\quad
       f\longmapsto\sum_{x\in H\backslash G}x^{-1}f(x).
                                                               \tag{S30.1}
\]
The inverse sum is independent of coset representatives. The summand
indexed by \(Hx\) is in \(x^{-1}W\); these are exactly the
distinct blocks. Thus both compositions are the identity.
Moreover, (S30.1) is an isometry for the form on V and the sum of the
forms on the values \(f(x)\).

\paragraph{Extension of Shapiro's map at every branch point.}
Remove the branch points of \(C\to\mathbb P^1\), obtaining
\(U\), and let \(Y_U=h^{-1}U\). Let \(\mathcal V\)
and \(\mathcal W\) be the finite local systems associated to
V on U and W on \(Y_U\). Over an unramified simply connected
disc, (S30.1) gives an isomorphism
\(\mathcal V\simeq h_*\mathcal W\).
Its monodromy compatibility is the identity
\(f_{gv}(x)=f_v(xg)\), so these disc maps glue.

At a removed base point, take a block cycle of length a for its local
monodromy g. A fixed vector in the corresponding induced summand is
determined by its value w in the first block, with successive values
obtained by the powers of g and with \(g^a w=w\). The same
operator \(g^a\), restricted to that block and conjugated by its
coset representative, is the actual local monodromy of q at the
corresponding point of Y. Thus the fixed stalks agree, including when
this coefficient monodromy is the identity. This gives an isomorphism
on the compact curves
\begin{equation}\label{eq:S30-compact-Shapiro}
       j_*\mathcal V\ \simeq\ h_*j'_*\mathcal W.
\end{equation}
Here \(j'\) fills every point at which the coefficient local
system is unramified; its remaining punctures are exactly the true
branch divisor D of q. A finite map has no higher direct images for
these constructible coefficient sheaves: over a small disc its inverse
image is a finite union of discs, and a finite-monodromy local system
extended by its invariant stalk has no positive cohomology there.
One can check the latter directly by passing to the finite cyclic
disc cover and taking the exact invariant summand. Consequently
\eqref{eq:S30-compact-Shapiro} gives
\begin{equation}\label{eq:S30-parabolic-Shapiro}
 H^1(\mathbb P^1,j_*\mathcal V)
            \simeq H^1(Y,j'_*\mathcal W).
\end{equation}
This is the explicit parabolic Shapiro map, including all original
branch points that disappear from D.

It preserves the Hermitian cup pairing. On the punctured curves,
integration over the inverse-image sheets is the sum of the integrals
on those sheets, and (S30.1) is an isometry for exactly that sum form.
The parabolic classes extend over the finite branch discs by the
invariant-stalk calculation, so filling those discs gives the same
pairing. There is no unmentioned factor N in this identity.
It also preserves the Hodge filtration. More explicitly, the map
\(C\to X\) gives actual pullback on differentials and cohomology;
its trace satisfies \(p_*p^*=[H_0:1]\). The components of
(S30.1) are block projections and deck operators, and averaging over
\(H_0\) is \(|H_0|^{-1}p^*p_*\). These pullback, trace,
and deck maps are the above local-system maps on the finite cover.
For an explicit representation-theoretic version, restriction sends a
G-equivariant coefficient map F to \(F|_W\), whose image is
\(H_0\)-fixed. Its inverse from an H-equivariant map f is
\[
       F(v)=\sum_{x\in H\backslash G}
                T_{x^{-1}}f\bigl(\pi_W\rho(x)v\bigr).
\]
Passing between the \(H_0\)-fixed image and \(H^1(X)\)
uses exactly \(p^*\) and \(|H_0|^{-1}p_*\). The same
formulas apply to the dual coefficient representations. Thus every
map used in Frobenius reciprocity here is a fixed sum of actual deck,
pullback and trace maps. They are morphisms of Hodge structures and
vary as fixed flat maps in a local topological marking.

\paragraph{Comparison with the actual polarized Jacobian factor.}
The primitive factor B of \(JC\) is obtained by the rational
rank-one matrix-coefficient idempotent on V together with its
conjugate, with denominators cleared, exactly as in HD-S23.
For clarity, take a vector v and divide its coefficient map by
\(h_V(v,v)\): on a regular cover fibre it is
\[
                 j_{v,x}(g)
           =\frac{h_V(v,\rho(g)x)}{h_V(v,v)}.
\]
The finite Schur identity gives the precise positive scalar
\[
 \sum_{g\in G}\overline{j_{v,x}(g)}j_{v,y}(g)
       =\frac{|G|}{\dim_K V\,h_V(v,v)}\,h_V(x,y).
                                                               \tag{S30.2}
\]
Write \(S(x,y)\) for the skew-Hermitian parabolic cup pairing,
conjugate-linear in x and linear in y, and choose the real alternating
convention \(E_{\rm par}(x,y)=\operatorname{Re}S(x,y)\).
The map from its underlying real space to B-homology is Poincare
duality applied to \(j_{v,x}+\overline{j_{v,x}}\); its K-action
is the actual central K-action on this idempotent factor. The
construction using the covariant Jacobian deck action and its
coefficient action is the one made explicit in HD-S23. The cup of two
unconjugated coefficient classes is zero, because their irreducible
complex block is not self-dual; its dual is the distinct conjugate
block. The two cross terms and (S30.2) therefore give exactly
\[
 E_B(\Phi x,\Phi y)
  =\frac{|G|}{\dim_KV\,h_V(v,v)}
           \bigl(S(x,y)+\overline{S(x,y)}\bigr)
  =\frac{2|G|}{\dim_KV\,h_V(v,v)}E_{\rm par}(x,y).
                                                               \tag{S30.3}
\]
In particular the normalization \(h_V(v,v)=2\) gives the
factor \(|G|/\dim_KV\), exactly as in HD-S23. This distinguishes
the complex Hermitian coefficient identity (S30.2) from the real
alternating map (S30.3). The rational idempotent, the denominator, and
this positive polarization multiplier are constant in the marked family.

Thus \eqref{eq:S30-parabolic-Shapiro} identifies the polarized
period of B with a fixed linear image of the coefficient variation
of \(X\to Y\). The integral image lattice may differ, but the
map is the specified fixed rational Hodge isomorphism from the actual
Jacobian morphisms above. A specified isogeny replaces it by another
fixed invertible rational map and has an inverse on period domains;
therefore it leaves the differential-rank bound unchanged.

\paragraph{Exact dimension and the low-genus moduli map.}
Put \(g=g(Y)\), \(r=|D|\), and let
\(c_a=\operatorname{rank}(1-u_a)\) for the nonidentity
coefficient inertia \(u_a\in\Gamma\) at \(a\in D\).
Every such codimension is an integer in \(\{1,2\}\).
The invariant-stalk Euler calculation gives
\[
 \chi(Y,j'_*\mathcal W)
  =2(2-2g-r)+\sum_{a\in D}(2-c_a)
  =4-4g-\sum c_a.
\]
The degree-zero cohomology is zero by irreducibility. The degree-two
cohomology is also zero: \(j'_*\mathcal W\) is the W-isotypic
coefficient summand of the actual finite map \(q_*\mathbb C_X\),
whereas \(H^2(X,\mathbb C)\) is the trivial deck representation
and W is nontrivial. Hence the six-dimensional parabolic space of B
satisfies the exact equality
\begin{equation}\label{eq:S30-dimension}
                   6=4g-4+\sum_{a\in D}c_a.
\end{equation}
It forces \(g\le2\). The possible ranges are
\[
\begin{array}{c|ccc}
g&0&1&2\\\hline
\sum c_a&10&6&2\\
r&5\le r\le10&3\le r\le6&1\le r\le2\\
3g-3+r&\le7&\le6&\le5.
\end{array}
\]
In every case \(2g-2+r>0\), so the actual pointed base curve
\((Y,D)\) is stable and has only finitely many automorphisms.
There is therefore no genus-zero or genus-one automorphism correction
missing from the displayed dimension.

Here is the family map giving the upper bound, rather than merely a
dimension comparison. Send the original G-cover to the quotient
\(X\to(Y,D)\) and discard the auxiliary map \(h:Y\to
\mathbb P^1\) and all coefficient-unramified marks. After a finite
base change one may order D and fix a marking of \(\Gamma\).
Over a small deformation of a stable ordered pointed curve, the
finite set of homomorphisms from its punctured topological fundamental
group to the fixed finite group \(\Gamma\), with the prescribed
inertia, is discrete. Retaining these permutations on a cut surface
constructs the unique continued unramified cover, and filling by the
fixed maps \(z\mapsto z^e\) gives its compact branched cover.
Thus the local coefficient-cover stack is a finite cover of the
pointed-curve deformation stack; the deck and base automorphisms are
finite. Its tangent dimension is at most \(3g-3+r\).
Equivalently this dimension is \(h^1(T_Y(-D))\): stability
kills \(h^0\), the curve has no degree-two deformation obstruction,
and Riemann--Roch gives the stated number.

By the fixed polarized map already proved, the original sixfold
period factors through this actual coefficient-cover map. Its
differential rank is at most \(3g-3+r\le7\), proving the
assertion.

\paragraph{The adjacent rank-three reduction.}
The same explicit construction for blocks of dimension three gives
\(6=6g-6+\sum c_a\), with \(1\le c_a\le3\).
At \(g=0,1,2\) the pointed coefficient-cover dimensions are at
most \(9,6,3\), respectively; larger g is impossible. Therefore
a period rank of nine requires exactly \(g=0\), twelve true
coefficient branch points, and \(c_a=1\) at every one.
The quotient coefficient curve is then \(\mathbb P^1\), and
its local monodromies generate the actual finite image \(\Gamma\).
They are all complex reflections on W. Thus a rank-nine source with
such three-dimensional blocks maps, by the displayed quotient and
Shapiro morphisms, to a twelve-reflection sphere-cover source in
rank three. No classification of the resulting reflection group is
needed or asserted here.

\paragraph{Proof audit.}
The complete author proof was read after writing. The maps of actual
quotient curves, the inverse induced-module formulas, the branch-stalk
extension, the positive cup scalar, and the stable pointed-cover
dimension are all supplied above. The rank bound uses the proved
factorization; it is not inferred from equal dimensions. There is no
new finite computation needed for the universal statement. The named
\(G_{12}\) instance uses the separately certified matrices of S29.
Independent review is recorded separately when it occurs.
